\scanpage{101}\mspage{- 171 -}\hfill IX

engendre $R$ au voisinage de $s$, d'où il résulte que $F(R)\subset R'$ au voisinage de $s$, ce qui achève la démonstration.

\uline{Corollaire} 13.2.2. \uline{Si} $A$ \uline{satisfait les conditions} a), b), c) \uline{de} 13.2.2, \uline{il est déterminé à isomorphisme unique près au voisinage de} $s$ \uline{par la connaissance du} $\hat S$-\uline{groupe} $\hat A$.

13.3. \uline{L'extension de \unsure{Raynaud} analytique}. Nous sommes maintenant en mesure de paraphraser, en Géométrie Analytique, la construction de l'extension de Raynaud donnée dans 7 . Supposons donné $A$ comme dans 13.1, avec $A=A_{\unsure{o|0}}$, et considérons $R_o=R_s$, qui est un $\mathbb Z$-module libre de type fini. Quitte à restreindre $S$ à un voisinage ouvert assez petit de $s$, interprétant $R_o$ comme le groupe des germes de sections de $R$ au voisinage de $s$, on trouve un homomorphisme
% Short unattached horizontal stroke above the space in « trouve un ».
\[
 R_{oS}\longrightarrow R\subset E \tag*{(13.3.1)}
\]
où $R_{oS}$ désigne le groupe analytique constant sur $S$ de valeur $R_o$. Choisissant le voisinage de $s$ assez petit pour que toute section de $\underline O_S$ dont le germe en $s$ est nul soit nul, on voit que l'homomorphisme précédent est injectif. Son image est donc un sous-groupe analytique constant de $R$, qui est évidemment le plus grand sous-groupe constant de $R$, et qui s'appelle la partie \uline{fixe} de $R$, et est notée $R^f$ :
% Short unattached horizontal stroke below the space in « un sous-groupe ».
\[
 R_{oS}\simeq R^f\lhook\joinrel\relbar R\,,\ \text{induit un isomorphisme } R^f_s\simeq R_s\,.
 \tag*{(13.3.2)}
\]
On pourra alors considérer le groupe quotient
\[
 A^{\handwritten{\natural}}=E/R^f\tag*{(13.3.3)}
\]
% archive sheet number: 100
% stamp: numbered label 489; faint pencil numerals above the lower strip.
